% GENERATED FILE. Edit canonical lessons, then run npm run generate:books.
% canonical-source-hash: fnv1a64:4c4636b5f5d0accf
% canonical-lessons: AR-C515-shakhara, AR-C515-istarkha, AR-C515-istalqa, AR-C515-nahada, AR-C515-ighma

\chapter{Lying down and getting up}
\label{ch:ar-lying-down-and-getting-up}
\hlchaptermodality{515}

\begin{chapteropening}
\textbf{By the end of this chapter:} \emph{I can say that someone snored, relaxed, lay down or got up, and talk about fainting.}

Complete the last lesson of chapter 515: I can say that someone snored, relaxed, lay down or got up, and talk about fainting.
\end{chapteropening}

\section[\texorpdfstring{\ar{شخر} (shakhara)}{شخر (shakhara)}]{\ar{شخر} (shakhara) — to snore}

\label{lesson:AR-C515-shakhara}



\begin{quote}
Before the new one: say the Arabic for to yawn, then the Arabic for to doze off.
\end{quote}

\subsection*{You'll want to know: \ar{شخر}}
\textbf{\ar{شخر}} — \emph{shakhara} — \textquotedblleft{}to snore\textquotedblright{}.

\begin{quote}
\ar{يشخر} \ar{أخي} \ar{في} \ar{الليل}.
\end{quote}

\begin{quote}
\emph{yashkhiru akhī fī al-layl.}
\end{quote}

\begin{quote}
My brother snores at night.
\end{quote}

Present: \ar{يشخر}, \emph{yashkhir}.

\begin{cousinweb}
\textbf{\ar{شخر}} (\emph{shakhara}) is built on the native Arabic root \emph{sh-kh-r}, \textquotedblleft{}to snort\textquotedblright{} (Lane). Lane gives shakhara, aor. yashkhiru, as a snorting sound from the nose or throat.
\end{cousinweb}

\begin{grammarlens}[title={What you've built}]
One more word for describing what you lived and felt.
\end{grammarlens}

\subsection*{Your turn}
\begin{itemize}
\raggedright
  \item \emph{Say it:} \emph{shakhara}
  \item \emph{Say it:} \emph{shakhara}, once more
  \item \emph{Say it:} \emph{yashkhiru akhī fī al-layl.} — My brother snores at night.
  \item \emph{Recall:} say the Arabic for to yawn, then the Arabic for to doze off, then say \emph{shakhara} again
\end{itemize}

\begin{tcolorbox}[breakable,colback=teal!4,colframe=teal!35!black,title={Before you move on}]
What does \ar{شخر} mean? (\textquotedblleft{}To snore\textquotedblright{}.) Say it once more.
\end{tcolorbox}
\section[\texorpdfstring{\ar{استرخى} (istarkhā)}{استرخى (istarkhā)}]{\ar{استرخى} (istarkhā) — to relax}

\label{lesson:AR-C515-istarkha}



\begin{quote}
Before the new one: say the Arabic for to doze off, then the Arabic for to snore.
\end{quote}

\subsection*{You'll want to know: \ar{استرخى}}
\textbf{\ar{استرخى}} — \emph{istarkhā} — \textquotedblleft{}to relax\textquotedblright{}.

\begin{quote}
\ar{استرخت} \ar{أمي} \ar{على} \ar{الأريكة}.
\end{quote}

\begin{quote}
\emph{istarkhat ummī ʿalā al-arīka.}
\end{quote}

\begin{quote}
My mother relaxed on the sofa.
\end{quote}

Present: \ar{يسترخي}, \emph{yastarkhī}.

\begin{cousinweb}
\textbf{\ar{استرخى}} (\emph{istarkhā}) is built on the native Arabic root \emph{r-kh-w}, \textquotedblleft{}to be slack, loose\textquotedblright{} (Lane). Form X istarkhā is to let the body go loose.
\end{cousinweb}

\begin{grammarlens}[title={What you've built}]
One more word for describing what you lived and felt.
\end{grammarlens}

\subsection*{Your turn}
\begin{itemize}
\raggedright
  \item \emph{Say it:} \emph{istarkhā}
  \item \emph{Say it:} \emph{istarkhā}, once more
  \item \emph{Say it:} \emph{istarkhat ummī ʿalā al-arīka.} — My mother relaxed on the sofa.
  \item \emph{Recall:} say the Arabic for to doze off, then the Arabic for to snore, then say \emph{istarkhā} again
\end{itemize}

\begin{tcolorbox}[breakable,colback=teal!4,colframe=teal!35!black,title={Before you move on}]
What does \ar{استرخى} mean? (\textquotedblleft{}To relax\textquotedblright{}.) Say it once more.
\end{tcolorbox}
\section[\texorpdfstring{\ar{استلقى} (istalqā)}{استلقى (istalqā)}]{\ar{استلقى} (istalqā) — to lie down}

\label{lesson:AR-C515-istalqa}



\begin{quote}
Before the new one: say the Arabic for to snore, then the Arabic for to relax.
\end{quote}

\subsection*{You'll want to know: \ar{استلقى}}
\textbf{\ar{استلقى}} — \emph{istalqā} — \textquotedblleft{}to lie down\textquotedblright{}.

\begin{quote}
\ar{استلقى} \ar{الولد} \ar{على} \ar{العشب}.
\end{quote}

\begin{quote}
\emph{istalqā al-walad ʿalā al-ʿushb.}
\end{quote}

\begin{quote}
The boy lay down on the grass.
\end{quote}

Present: \ar{يستلقي}, \emph{yastalqī}.

\begin{cousinweb}
\textbf{\ar{استلقى}} (\emph{istalqā}) is built on the native Arabic root \emph{l-q-y}, \textquotedblleft{}to meet; to throw\textquotedblright{} (Lane). Form X istalqā is to throw oneself down on one's back.
\end{cousinweb}

\begin{grammarlens}[title={What you've built}]
One more word for describing what you lived and felt.
\end{grammarlens}

\subsection*{Your turn}
\begin{itemize}
\raggedright
  \item \emph{Say it:} \emph{istalqā}
  \item \emph{Say it:} \emph{istalqā}, once more
  \item \emph{Say it:} \emph{istalqā al-walad ʿalā al-ʿushb.} — The boy lay down on the grass.
  \item \emph{Recall:} say the Arabic for to snore, then the Arabic for to relax, then say \emph{istalqā} again
\end{itemize}

\begin{tcolorbox}[breakable,colback=teal!4,colframe=teal!35!black,title={Before you move on}]
What does \ar{استلقى} mean? (\textquotedblleft{}To lie down\textquotedblright{}.) Say it once more.
\end{tcolorbox}
\section[\texorpdfstring{\ar{نهض} (nahaḍa)}{نهض (nahaḍa)}]{\ar{نهض} (nahaḍa) — to get up, to rise}

\label{lesson:AR-C515-nahada}



\begin{quote}
Before the new one: say the Arabic for to relax, then the Arabic for to lie down.
\end{quote}

\subsection*{You'll want to know: \ar{نهض}}
\textbf{\ar{نهض}} — \emph{nahaḍa} — \textquotedblleft{}to get up, to rise\textquotedblright{}.

\begin{quote}
\ar{نهض} \ar{جدي} \ar{من} \ar{الكرسي} \ar{ببطء}.
\end{quote}

\begin{quote}
\emph{nahaḍa jaddī min al-kursī bi-buṭʾ.}
\end{quote}

\begin{quote}
My grandfather got up slowly from the chair.
\end{quote}

Present: \ar{ينهض}, \emph{yanhaḍ}.

The noun \ar{نهضة} is a rising, a revival.

\begin{cousinweb}
\textbf{\ar{نهض}} (\emph{nahaḍa}) is built on the native Arabic root \emph{n-h-ḍ}, \textquotedblleft{}to rise, to stand up\textquotedblright{} (Lane). Nahaḍa is to rise up from sitting.
\end{cousinweb}

\begin{grammarlens}[title={What you've built}]
One more word for describing what you lived and felt.
\end{grammarlens}

\subsection*{Your turn}
\begin{itemize}
\raggedright
  \item \emph{Say it:} \emph{nahaḍa}
  \item \emph{Say it:} \emph{nahaḍa}, once more
  \item \emph{Say it:} \emph{nahaḍa jaddī min al-kursī bi-buṭʾ.} — My grandfather got up slowly from the chair.
  \item \emph{Recall:} say the Arabic for to relax, then the Arabic for to lie down, then say \emph{nahaḍa} again
\end{itemize}

\begin{tcolorbox}[breakable,colback=teal!4,colframe=teal!35!black,title={Before you move on}]
What does \ar{نهض} mean? (\textquotedblleft{}To get up, to rise\textquotedblright{}.) Say it once more.
\end{tcolorbox}
\section[\texorpdfstring{\ar{إغماء} (ighmāʾ)}{إغماء (ighmāʾ)}]{\ar{إغماء} (ighmāʾ) — fainting, a faint}

\label{lesson:AR-C515-ighma}



\begin{quote}
Before the new one: say the Arabic for to lie down, then the Arabic for to get up.
\end{quote}

\subsection*{You'll want to know: \ar{إغماء}}
\textbf{\ar{إغماء}} — \emph{ighmāʾ} — \textquotedblleft{}fainting, a faint\textquotedblright{}.

\begin{quote}
\ar{نقلوا} \ar{اللاعب} \ar{إلى} \ar{المستشفى} \ar{بعد} \ar{الإغماء}.
\end{quote}

\begin{quote}
\emph{naqalū al-lāʿib ilā al-mustashfā baʿda al-ighmāʾ.}
\end{quote}

\begin{quote}
They took the player to hospital after he fainted.
\end{quote}

\begin{cousinweb}
\textbf{\ar{إغماء}} (\emph{ighmāʾ}) is built on the native Arabic root \emph{gh-m-y}, \textquotedblleft{}to cover, to veil\textquotedblright{} (Lane). Ighmāʾ is the verbal noun of Form IV: being covered over, losing one's senses.
\end{cousinweb}

\begin{grammarlens}[title={What you've built}]
One more word for describing what you lived and felt.
\end{grammarlens}

\subsection*{Your turn}
\begin{itemize}
\raggedright
  \item \emph{Say it:} \emph{ighmāʾ}
  \item \emph{Say it:} \emph{ighmāʾ}, once more
  \item \emph{Say it:} \emph{naqalū al-lāʿib ilā al-mustashfā baʿda al-ighmāʾ.} — They took the player to hospital after he fainted.
  \item \emph{Recall:} say the Arabic for to lie down, then the Arabic for to get up, then say \emph{ighmāʾ} again
\end{itemize}

\begin{tcolorbox}[breakable,colback=teal!4,colframe=teal!35!black,title={Before you move on}]
What does \ar{إغماء} mean? (\textquotedblleft{}Fainting, a faint\textquotedblright{}.) Say it once more.
\end{tcolorbox}
